\documentclass[11pt,a4paper]{article}

\usepackage[T1]{fontenc}
\usepackage[utf8]{inputenc}
\usepackage{amsmath,amssymb,amsthm,mathtools}
\usepackage[margin=1.1in]{geometry}
\usepackage[expansion=false]{microtype}
\usepackage{booktabs}
\usepackage{longtable}
\usepackage[colorlinks=true,linkcolor=blue,citecolor=blue,urlcolor=blue]{hyperref}

\theoremstyle{plain}
\newtheorem{theorem}{Theorem}[section]
\newtheorem{proposition}[theorem]{Proposition}
\newtheorem{lemma}[theorem]{Lemma}
\newtheorem{corollary}[theorem]{Corollary}
\theoremstyle{definition}
\newtheorem{definition}[theorem]{Definition}
\newtheorem{remark}[theorem]{Remark}
\newtheorem{question}[theorem]{Question}

\DeclareMathOperator{\conv}{conv}
\DeclareMathOperator{\sgn}{sgn}
\DeclareMathOperator{\vertt}{vert}
\DeclareMathOperator{\supp}{supp}
\DeclareMathOperator{\VCn}{\mathrm{VC}_n}
\newcommand{\R}{\mathbb{R}}
\newcommand{\N}{\mathbb{N}}
\newcommand{\Ftwo}{\mathbb{F}_2}
\newcommand{\Grid}{\mathcal{G}}
\newcommand{\lean}[1]{\texttt{\footnotesize #1}}

\title{$\VCn$ dimension of convex sets in $\R^m$}

\author{LLMs prompted by Yaël Dillies}
\date{\today}

\begin{document}
\maketitle

\begin{abstract}
For a subset $A$ of an abelian group $G$, the \emph{$\VCn$ dimension} of $A$
is the maximal natural number $d$ such that there exist finite families $X_1, \dots, X_n$
of $d$ elements such that $A$ shatters $X_1 + \dots + X_n$.
We study the maximum $\VCn$ dimension of a convex set $C \subseteq \R^m$.
For $m = n$, we prove an upper bound of $2$, which is tight.
For $m = n + 2$, we construct a convex set $C$ of infinite $\VCn$ dimension,
generalising the construction of Naszódi and Taschuk \cite[Example~4]{naszoditaschuk}.
For $m = n + 1$, we prove an upper bound of $2^{2^{n + 1} + 2} - 1$ by reducing to bounding
the number of sign patterns of the minors of an additive grid following Warren \cite{warren},
and separately bounding the number of convex-independent subsets of an additive grid.
\end{abstract}

\tableofcontents

\section{Introduction}\label{sec:intro}

\subsection{The problem}

Let $G$ be an abelian group and $A\subseteq G$. Fix $n,d\in\N$ and write
$[m]:=\{0,\dots,m-1\}$.

\begin{definition}[$\VCn$ dimension]\label{def:vcn}
A set $A\subseteq G$ \emph{has $\VCn$ dimension at most $d$}, written
$\VCn(A)\le d$, if there is \emph{no} pair
\[
  x\colon [n]\times[d+1]\to G,
  \qquad
  y\colon \mathcal{P}\bigl([n]\to[d+1]\bigr)\to G
\]
such that
\begin{equation}\label{eq:shatter}
  y(S)+\sum_{k<n} x\bigl(k,i(k)\bigr)\in A
  \iff
  i\in S
  \qquad\text{for all } i\colon[n]\to[d+1] \text{ and all } S.
\end{equation}
\end{definition}

In words: $\VCn(A)>d$ exactly when there is an $n$-dimensional grid of
``directions'' $x$ of side length $d+1$ whose $(d+1)^n$ sums can be translated so as to
realise \emph{every} subset of the grid inside $A$. For $n=1$ this is the ordinary VC
dimension of the family of translates of $A$.

\subsection{Origins of the \texorpdfstring{$\VCn$}{VCn} dimension}

The case $n=1$ goes back to Vapnik and Chervonenkis \cite{vc}, who introduced the dimension now
named after them to prove uniform laws of large numbers. At the same time, and independently,
Shelah \cite{shelah71} isolated the \emph{independence property} of a formula in model theory;
Laskowski \cite{laskowski} later observed that a formula is dependent (NIP) exactly when the
family of sets it defines has finite VC dimension.

The higher-arity notion also comes from model theory. Shelah \cite{shelahstrongly,shelah2dep}
introduced \emph{$n$-dependent} theories, in which no formula $\varphi(x_1,\dots,x_n;y)$ can
encode a random $(n+1)$-partite $(n+1)$-uniform hypergraph. Chernikov, Palacín and Takeuchi
\cite{cpt} developed the combinatorics of this notion. In their terms, an $(n+1)$-ary relation
$R\subseteq X_1\times\dots\times X_n\times Y$ has $\VCn$ dimension at least $d+1$ if there
are $a_{k,j}\in X_k$ ($k<n$, $j\le d$) and $b_S\in Y$ ($S\subseteq[d+1]^{n}$) with
$R(a_{0,i(0)},\dots,a_{n-1,i(n-1)};b_S)$ if and only if $i\in S$. They showed that a formula is
$n$-dependent exactly when it has finite $\VCn$ dimension, and proved a generalisation of
the Sauer--Shelah lemma for it. Chernikov and Towsner \cite{chernikovtowsner} then showed that
hypergraphs of bounded $\VCn$ dimension admit hypergraph regularity decompositions that
use only the first $n$ levels.

Definition~\ref{def:vcn} is the specialisation to the relation $y+x_1+\dots+x_n\in A$
on an abelian group. For $n=1$, sets of bounded VC dimension in abelian groups have strong
arithmetic regularity properties, studied by Alon, Fox and Zhao \cite{afz} and by Sisask
\cite{sisask}. For $n=2$, Terry and Wolf \cite{terrywolf,terrywolfvc2} proved that subsets of
$\mathbb{F}_p^{m}$ of bounded $\mathrm{VC}_2$ dimension are well approximated by unions of atoms
of a bounded-complexity quadratic factor. Gladkova \cite{gladkova} confirmed their conjecture that
the ``quadratic Green--Sanders example'' in $\mathbb{F}_p^{m}$ has $\mathrm{VC}_2$ dimension at
least $3$ (and at most $501$). Convex sets are a natural Euclidean test case: in $\R^{n+1}$ the
$\VCn$ dimension is finite (Theorem~\ref{thm:main}), and in $\R^{n+2}$ it need not be
(Theorem~\ref{thm:infinite}).

\begin{definition}\label{def:dstar}
For $n\ge1$ let $d^\ast(n)$ be the least $d\in\N\cup\{\infty\}$ such that
$\VCn(C)\le d$ for every convex $C\subseteq\R^{n+1}$.
\end{definition}

\begin{question}[The problem]\label{q:main}
Determine $d^\ast(n)$. In particular:
\begin{enumerate}
\item[(a)] Is $d^\ast(n)$ finite for every $n\ge1$?
\item[(b)] Is $d^\ast(n)\le3$ for every $n\ge1$?
\item[(c)] Is $d^\ast(2)\le2$, that is, does every convex $C\subseteq\R^3$ satisfy
$\mathrm{VC}_2(C)\le2$?
\end{enumerate}
\end{question}

Question~(a) is the qualitative problem and is answered in this paper. Questions (b) and (c) ask
for the true value: a uniform bound independent of $n$, and for $n=2$ a bound one smaller
still. For $n=1$ the value is at most $3$ by Naszódi and Taschuk (Theorem~\ref{thm:planar}). Whether $d^\ast(n)$ grows with $n$, and
at what rate, is open; the bound proved here is doubly exponential in $n$.

\subsection{The result}

\begin{theorem}\label{thm:main}
Let $n\ge1$ and put
\[
  d(n)\;:=\;2^{2^{n+1}+2}-1 .
\]
Then every convex $C\subseteq\R^{n+1}$ satisfies $\VCn(C)\le d(n)$. In
particular $d^\ast(n)\le d(n)<\infty$, which answers Question~\ref{q:main}(a).
\end{theorem}

In small dimensions much better bounds hold. The first is known.

\begin{theorem}[Naszódi--Taschuk {\cite[Theorem~3]{naszoditaschuk}}]\label{thm:planar}
Translates of a convex set $C\subseteq\R^{2}$ do not shatter any set of four points. In
particular $\mathrm{VC}_1(C)\le3$, so $d^\ast(1)\le3$, which is Question~\ref{q:main}(b) for $n=1$.
\end{theorem}

\begin{theorem}\label{thm:three}
Every convex $C\subseteq\R^{3}$ satisfies $\mathrm{VC}_2(C)\le123$.
\end{theorem}

The dimension $n+1$ in Theorem~\ref{thm:main} cannot be increased. For $n=1$ this is the
example of Naszódi and Taschuk \cite[Example~4]{naszoditaschuk} of a convex body in $\R^{3}$
whose translates have infinite VC dimension, and the following theorem extends it to all $n$.

\begin{theorem}\label{thm:infinite}
For every $n\ge0$ there is a convex set $C\subseteq\R^{n+2}$ with $\VCn(C)=\infty$, i.e.\
$\VCn(C)\le d$ fails for every $d$.
\end{theorem}

For $n\ge2$ Questions~\ref{q:main}(b) and (c) remain open: the bound $d(n)$ is doubly
exponential, whereas the conjectured truth is $3$.

\subsection{Strategy}

For $n\ge2$, everything rests on a counting reduction. Fix $D:=n+1$ and a side length $m$. An
\emph{additive array} is a map $z\colon[D]\times[m]\to\R^D$, and it assigns to each grid
index $i\in\Grid:=[m]^{[D]}$ the point
\[
  \Sigma z(i)\;:=\;\sum_{k<D}z\bigl(k,i(k)\bigr)\in\R^D .
\]
Call $S\subseteq\Grid$ a \emph{convex label} if there is a convex $C\subseteq\R^D$ and an
additive array $z$ with $S=\{i:\Sigma z(i)\in C\}$, and let $\mathcal{L}(D,m)$ be the set
of convex labels.

The whole proof is the inequality
\begin{equation}\label{eq:target}
  \bigl|\mathcal{L}(D,m)\bigr| \;<\; 2^{m^{D}}
  \qquad\text{for}\qquad m=2^{2^{D}+2},
\end{equation}
that is: \emph{not every} subset of the grid is a convex label. Since shattering at scale
$m$ would force \emph{all} $2^{m^D}$ subsets to be labels (Proposition~\ref{prop:shatter}),
\eqref{eq:target} gives $\mathrm{VC}_{D-1}(C)\le m-1$ for every convex $C$, which is
Theorem~\ref{thm:main}.

Inequality~\eqref{eq:target} is proved by exhibiting an injective code for convex labels
with two coordinates:
\[
  \mathcal{L}(D,m)\;\hookrightarrow\;
  \underbrace{\Sigma(D,m)}_{\text{sign vectors of all minors}}\times
  \underbrace{\mathcal{G}_K}_{\text{subsets of size}\le K} .
\]
The first factor is the classical polynomial-method coordinate: a label is determined by
the signs of all the determinants built from the homogenised array, and sign patterns of
low-degree polynomials are few (Section~\ref{sec:signs}). The second factor requires a
genuinely convex-geometric input: the set of hull generators of a label is convexly independent,
and a convexly independent set of indices has at most $K=V(n,m)$ elements, which is at most a
quarter of the grid at our choice of $m$ (Section~\ref{sec:vertices}). Section~\ref{sec:numerics}
carries out the arithmetic. For $n=1$ this arithmetic fails, and Section~\ref{sec:planar}
deduces the case $n=1$ from Theorem~\ref{thm:planar} instead. Section~\ref{sec:three} sketches
Theorem~\ref{thm:three}, and Section~\ref{sec:infinite} proves Theorem~\ref{thm:infinite}.
Section~\ref{sec:status} discusses the gap to the conjectured bounds,
and Section~\ref{sec:formal} describes a formalisation of all the results in Lean~4.

The main new ingredient is Theorem~\ref{thm:fibre}, a bound on convexly independent subsets of
an additive grid, proved by an induction on the number of colours which carries along a subspace of
directions annihilated by every exposing functional.

\subsection{Conventions}

Throughout, $D\ge1$ is an integer (so $n=D-1$), $m\ge1$ is a side length,
$\Grid=\Grid(D,m)=[m]^{[D]}$ is the grid of size $M:=m^D$. For
$x\in\R^D$ we write $\hat x:=(1,x)$ for its homogenisation, with the extra
coordinate indexed by a symbol $\ast$, so that rows of homogenised matrices are indexed by
$\{\ast\}\sqcup[D]$. All convex hulls are over $\R$. We use the three-valued sign code
$\sgn\colon\R\to\{-,0,+\}$.

\section{From shattering to a counting problem}\label{sec:counting}

\begin{proposition}[Shattering realises all labels]\label{prop:shatter}
Let $G$ be an abelian group, $C\subseteq G$, $n,m\ge1$, and suppose that
$x\colon[n]\times[m]\to G$ and $y$ witness \eqref{eq:shatter} for $C$ at side length
$m=d+1$. Then every $S\subseteq\Grid(n+1,m)$ is realised by an additive array in $C$; that
is, there is $z\colon[n+1]\times[m]\to G$ with $S=\{i:\Sigma z(i)\in C\}$.
\end{proposition}

\begin{proof}
Given $S$, and for $j\in[m]$, write $S_j:=\{i'\in\Grid(n,m) : (j,i')\in S\}$, where
$(j,i')$ denotes the tuple in $\Grid(n+1,m)$ with first entry $j$ and tail $i'$. Define
\[
  z(0,j):=y(S_j),\qquad z(k+1,\cdot):=x(k,\cdot)\quad(k<n).
\]
Then for $i=(j,i')$,
\[
  \Sigma z(i) \;=\; y(S_j)+\sum_{k<n}x\bigl(k,i'(k)\bigr),
\]
which lies in $C$ iff $i'\in S_j$ by \eqref{eq:shatter}, i.e.\ iff $i\in S$.
\end{proof}

The point of Proposition~\ref{prop:shatter} is that a shattering configuration uses only
$m$ of the $2^{m^n}$ available translations, but nevertheless forces \emph{every} label to
occur, \emph{with $C$ held fixed}. Consequently:

\begin{corollary}\label{cor:reduction}
Let $n\ge1$ and $m\ge1$. If $|\mathcal{L}(n+1,m)|<2^{m^{n+1}}$, then every convex
$C\subseteq\R^{n+1}$ satisfies $\VCn(C)\le m-1$.
\end{corollary}

\begin{proof}
If not, fix $C$ and a shattering pair $(x,y)$ at side length $m$. By
Proposition~\ref{prop:shatter} every $S\subseteq\Grid(n+1,m)$ is a convex label, so
$|\mathcal{L}(n+1,m)|=2^{m^{n+1}}$, a contradiction.
\end{proof}

From here on the ambient convex
set plays no further role: we must count labels.

\section{Coding a label}\label{sec:code}

\subsection{Reduction to a finite hull}

Fix an additive array $z$ and write $q:=\Sigma z\colon\Grid\to\R^D$ for the induced
configuration of $M=m^D$ points (possibly with repetitions).

\begin{lemma}[Hull encoding]\label{lem:hull}
Let $C\subseteq\R^D$ be convex and $q\colon\Grid\to\R^D$ arbitrary. Then there is
$V\subseteq\Grid$ such that
\begin{enumerate}
\item $q(i)\in C\iff q(i)\in\conv q(V)$ for all $i\in\Grid$; and
\item $V$ is \emph{convex independent}: $q(i)\notin\conv q(V\setminus\{i\})$ for all
      $i\in V$ (in particular $q|_V$ is injective).
\end{enumerate}
\end{lemma}

\begin{proof}
Let $E:=\{i: q(i)\in C\}$. Since $C$ is convex and $q(E)\subseteq C$, we get
$\conv q(E)\subseteq C$, so $q(i)\in C\iff q(i)\in\conv q(E)$. Choose $V\subseteq E$ of
minimal cardinality with $\conv q(V)=\conv q(E)$. If some $i\in V$ satisfied
$q(i)\in\conv q(V\setminus\{i\})$ then $\conv q(V\setminus\{i\})=\conv q(V)$, contradicting
minimality.
\end{proof}

Thus a convex label $S$ is recovered from the pair $(q,V)$ by
$S=\{i : q(i)\in\conv q(V)\}$. Neither closedness nor boundedness of $C$ is used, and $q$
may be degenerate (repeated points, lower-dimensional affine hull).

\subsection{Minors determine finite hull membership}

We now replace $q$ by a finite sign vector. Let
\[
  A(q)\;\in\;\R^{(\{\ast\}\sqcup[D])\times\Grid},
  \qquad
  A(q)_{\ast,i}=1,\quad A(q)_{a,i}=q(i)_a ,
\]
be the homogenised configuration matrix, whose $i$-th column is $\widehat{q(i)}$. A
\emph{minor index} is a pair $(R,c)$ with $R\subseteq\{\ast\}\sqcup[D]$ and
$c\colon R\to\Grid$ arbitrary (repetitions allowed); the associated minor is
$\det\bigl(A(q)_{R,c}\bigr)$. Say $q,q'$ have \emph{the same minor signs} if
$\sgn\det A(q)_{R,c}=\sgn\det A(q')_{R,c}$ for all $(R,c)$.

\begin{theorem}[Minor signs determine hulls]\label{thm:minorhull}
If $q,q'\colon\Grid\to\R^D$ have the same minor signs, then for every finite
$V\subseteq\Grid$ and every $i\in\Grid$,
\[
  q(i)\in\conv q(V)\iff q'(i)\in\conv q'(V).
\]
\end{theorem}

\begin{proof}
It suffices to prove one implication. Suppose $q(i)\in\conv q(V)$. Homogenising, convex
membership is the statement that $\widehat{q(i)}$ is a nonnegative combination of the
columns $\widehat{q(v)}$, $v\in V$ (the $\ast$-row enforces that the weights sum to $1$).
By Carathéodory's theorem there are $r\le D+1$, an injection
$c\colon[r]\to V$ with $\bigl(\widehat{q(c(j))}\bigr)_j$ linearly independent, and weights
$w_j>0$ with $\sum_j w_j\widehat{q(c(j))}=\widehat{q(i)}$.

Write $B:=A(q)$, $B':=A(q')$. Three consequences of equal minor signs:
\begin{enumerate}
\item[(a)] \emph{Independence transfers.} A family of columns of $B$ is independent iff some
square row restriction of it is nonsingular; the corresponding minor of $B'$ has the same
sign, hence is also nonzero, so the columns $\bigl(\widehat{q'(c(j))}\bigr)_j$ are
independent as well.
\item[(b)] \emph{Span membership transfers.} Applying (a) to the family $c$ enlarged by the
column $i$ shows that $B\,e_i$ lies in the span of the $B\,e_{c(j)}$ iff the same
holds for $B'$. Hence there are (a priori arbitrary) reals $v_j$ with
$\sum_j v_j\widehat{q'(c(j))}=\widehat{q'(i)}$.
\item[(c)] \emph{Signs of the weights transfer.} Choose rows $\rho\colon[r]\hookrightarrow
\{\ast\}\sqcup[D]$ making $B_{\rho,c}$ nonsingular; then $B'_{\rho,c}$ is nonsingular too.
On these rows, $w$ and $v$ are the unique solutions of the respective square systems, so by
Cramer's rule
\[
  w_j=\frac{\det\bigl(B_{\rho,c}\text{ with column }j\text{ replaced by }B_{\rho,i}\bigr)}
           {\det B_{\rho,c}},
\]
and likewise for $v_j$. Both numerator and denominator are minors of the family indexed by
$(\rho,\,c\text{ with }j\mapsto i)$ and $(\rho,c)$ respectively; their signs agree between
$q$ and $q'$. Hence $w_j\ge0$ implies $v_j\ge0$.
\end{enumerate}
So $v\ge0$, and $\sum_j v_j\widehat{q'(c(j))}=\widehat{q'(i)}$ exhibits
$q'(i)\in\conv q'(V)$.
\end{proof}

The proof deliberately avoids general position: the ambient configuration may be
rank-deficient, and the nonsingular square block is selected locally.

\subsection{The product bound}

Let $\Sigma(D,m)$ be the set of sign vectors $\bigl(\sgn\det A(\Sigma z)_{R,c}\bigr)_{(R,c)}$
realised by additive arrays $z$, and let $\mathcal{G}$ be any family of subsets of $\Grid$
which contains a witness $V$ as in Lemma~\ref{lem:hull} for every convex label.

\begin{corollary}\label{cor:product}
$\bigl|\mathcal{L}(D,m)\bigr|\le\bigl|\Sigma(D,m)\bigr|\cdot|\mathcal{G}|$.
\end{corollary}

\begin{proof}
Map each label $S$ to the pair (minor sign vector of its chosen array, its chosen $V$).
If two labels $S,T$ receive the same pair, then by Theorem~\ref{thm:minorhull} their
defining conditions $q(i)\in\conv q(V)$ and $q'(i)\in\conv q'(V)$ agree for all $i$, so
$S=T$.
\end{proof}

It remains to bound $|\Sigma(D,m)|$ (Section~\ref{sec:signs}) and to produce a small
family $\mathcal{G}$ (Section~\ref{sec:vertices}).

\section{Counting sign patterns}\label{sec:signs}

The classical tool here is Warren's theorem.

\begin{theorem}[Warren {\cite[Theorem~3]{warren}}]\label{thm:warren}
Let $f_1,\dots,f_N\in\R[x_1,\dots,x_p]$ with $\deg f_i\le k$ for all $i$, and $N\ge p$. Then the
number of sign words $s\in\{+,-\}^N$ for which there exists $x\in\R^p$ with
$\sgn f_i(x)=s_i$ for all $i$ is at most $(4ekN/p)^{p}$.
\end{theorem}

We need the version which also records zero signs. It follows by a standard perturbation.

\begin{corollary}[Ternary sign count]\label{cor:ternary}
With $f_i$ as above and $2N\ge p$, the number of realised \emph{ternary} sign vectors
$\bigl(\sgn f_i(x)\bigr)_{i}\in\{-,0,+\}^N$ is at most $(8ekN/p)^{p}$.
\end{corollary}

\begin{proof}
Choose a witness $x$ for each of the finitely many realised ternary vectors, and $\varepsilon>0$
smaller than $|f_i(x)|$ for every witness $x$ and every $i$ with $f_i(x)\ne0$. At a witness $x$,
\[
  \bigl(\sgn(f_i(x)-\varepsilon),\ \sgn(f_i(x)+\varepsilon)\bigr)=
  \begin{cases}
   (-,-) & f_i(x)<0,\\
   (-,+) & f_i(x)=0,\\
   (+,+) & f_i(x)>0,
  \end{cases}
\]
so the ternary vector is recovered injectively from a strict sign word of the $2N$ polynomials
$f_i\pm\varepsilon$. Apply Theorem~\ref{thm:warren} to them.
\end{proof}

Now specialise. The variables are the $D^2m$ coordinates of the array $z$; the entries of
the homogenised matrix are the polynomials
\[
  \ast\text{-row}:\ 1,
  \qquad
  a\text{-row}:\ \sum_{k<D}X_{(k,\,i(k),\,a)} \ \ (a<D),
\]
of degree $0$ and $1$ respectively. A minor with row set $R$ therefore has total degree at
most $\sum_{r\in R}\deg(\text{row } r)\le D$, \emph{including} the maximal minors of size
$D+1$. Allowing repeated columns (whose determinant is $0$) avoids a factorial, and the
number of minor indices is at most
\[
  \mathcal{N}(D,m)\;=\;2^{D+1}\bigl(m^{D}\bigr)^{D+1}.
\]

\begin{theorem}\label{thm:minorsigns}
For $D\ge1$ and $m\ge1$,
\[
  \bigl|\Sigma(D,m)\bigr|\;\le\;\Bigl(\frac{8e\,\mathcal{N}(D,m)}{Dm}\Bigr)^{D^2m}.
\]
\end{theorem}

\begin{proof}
Corollary~\ref{cor:ternary} with $k=D$, $N=\mathcal{N}(D,m)$ and $p=D^2m\le2N$.
\end{proof}

At a side length which is a power of two, the bound takes a simpler shape.

\begin{corollary}\label{cor:signpow}
Let $D\ge1$, $L\ge8$ and $m=2^{L}$. Then $\bigl|\Sigma(D,m)\bigr|\le2^{\,4L(D+1)^{4}m}$.
\end{corollary}

\begin{proof}
The base in Theorem~\ref{thm:minorsigns} is at most
$2^{5}\cdot2^{D+1}\cdot2^{LD(D+1)}\le2^{2L(D+1)^{2}}$, and the exponent satisfies
$D^{2}m\le2(D+1)^{2}m$.
\end{proof}

\section{Convexly independent subsets of an additive grid}\label{sec:vertices}

We now produce the second coordinate of the code. By Lemma~\ref{lem:hull} the generator set $V$
of a label can be taken convexly independent, so it suffices to show that convexly independent
sets of indices are small. The argument inducts on the number of colours, and uses nothing but
linear algebra and the Cauchy--Schwarz inequality.

\subsection{Exposing functionals}

For an array $a\colon[r]\times[m]\to\R^D$ and $x\in[m]^{[r]}$ write
$\Sigma a(x):=\sum_{k<r}a\bigl(k,x(k)\bigr)$. Say that a functional $\nu\in(\R^D)^{*}$
\emph{strictly exposes} $x$ in $S\subseteq[m]^{[r]}$ if $\nu(\Sigma a(y))<\nu(\Sigma a(x))$ for
all $y\in S\setminus\{x\}$.

\begin{lemma}\label{lem:expose}
If $V$ is convexly independent for $q$, then every $i\in V$ is strictly exposed in $V$ by some
functional: $\nu(q(j))<\nu(q(i))$ for all $j\in V\setminus\{i\}$.
\end{lemma}

\begin{proof}
The hull of the finitely many points $q(j)$, $j\ne i$, is compact and does not contain $q(i)$;
apply the Hahn--Banach separation theorem.
\end{proof}

\begin{lemma}\label{lem:combine}
Let $\nu,\nu'$ be linear functionals on a real vector space and $u$ a vector with $\nu(u)<0<\nu'(u)$.
Then $\nu_0:=\nu'(u)\,\nu-\nu(u)\,\nu'$ satisfies $\nu_0(u)=0$; every strict inequality
$\nu(x)<\nu(y)$ which also holds for $\nu'$ holds for $\nu_0$; and every vector annihilated by both
$\nu$ and $\nu'$ is annihilated by $\nu_0$.
\end{lemma}

\begin{proof}
$\nu_0$ is a combination of $\nu$ and $\nu'$ with positive coefficients $\nu'(u)$ and $-\nu(u)$.
\end{proof}

\begin{lemma}\label{lem:twopoints}
Let $W\le\R^D$ have codimension at most $1$. Then no three indices of a set $S$ are strictly
exposed in $S$ by functionals vanishing on $W$.
\end{lemma}

\begin{proof}
The functionals vanishing on $W$ form a space of dimension at most $1$, so they are the multiples
$c\,\varphi$ of a single $\varphi$. A multiple exposing a point strictly has $c\ne0$. Of three
exposed indices, two, say $x$ and $y$, are exposed by multiples with the same sign; but then
$\varphi(\Sigma a(y))<\varphi(\Sigma a(x))$ and $\varphi(\Sigma a(x))<\varphi(\Sigma a(y))$, or
the same with both inequalities reversed.
\end{proof}

\subsection{The fibre induction}

\begin{lemma}[Cauchy--Schwarz for intersecting families]\label{lem:cs}
Let $F_1,\dots,F_k$ be subsets of a finite set $X$ with $|F_i\cap F_j|\le K$ for all $i\ne j$.
Then
\[
  \Bigl(\sum_i|F_i|\Bigr)^{2}\;\le\;|X|\Bigl(\sum_i|F_i|+k(k-1)K\Bigr).
\]
\end{lemma}

\begin{proof}
Let $d(x):=|\{i:x\in F_i\}|$. Then $\sum_xd(x)=\sum_i|F_i|$ and
$\sum_xd(x)^2=\sum_{i,j}|F_i\cap F_j|\le\sum_i|F_i|+k(k-1)K$. Conclude by
$(\sum_xd(x))^2\le|X|\sum_xd(x)^2$.
\end{proof}

Define the \emph{vertex bound}
\[
  V(0,m):=2,\qquad
  V(r+1,m):=m^{r+1}+\Bigl\lfloor\sqrt{m^{r+1}\cdot m^{2}\cdot V(r,m)}\Bigr\rfloor .
\]

\begin{theorem}[Fibre induction]\label{thm:fibre}
Let $r,m,D\in\N$, let $a\colon[r+1]\times[m]\to\R^D$, and let $W\le\R^D$ be a subspace with
$D\le r+1+\dim W$. Let $S\subseteq[m]^{[r+1]}$ be such that every $x\in S$ is strictly exposed
in $S$ by a functional vanishing on $W$. Then $|S|\le V(r,m)$.
\end{theorem}

\begin{proof}
Induct on $r$. For $r=0$, $W$ has codimension at most $1$ and Lemma~\ref{lem:twopoints} gives
$|S|\le2$.

Let $r\ge0$, and consider $r+2$ colours. For $i\in[m]$ let
$F_i:=\{t\in[m]^{[r+1]}:(i,t)\in S\}$ be the fibre of $S$ over $i$, where $(i,t)$ denotes the index
with first entry $i$ and tail $t$. Then $\sum_i|F_i|=|S|$. Write $a'$ for the array $a$ with the
first colour removed, so that $\Sigma a(i,t)=a(0,i)+\Sigma a'(t)$.

\smallskip\noindent\emph{Two fibres share few tails.} Fix $i\ne i'$, put $T:=F_i\cap F_{i'}$ and
$u:=a(0,i')-a(0,i)$. For $t\in T$ let $\nu$ and $\nu'$ be the functionals exposing $(i,t)$ and
$(i',t)$. Comparing $(i,t)$ with $(i',t)$ gives $\nu(u)<0<\nu'(u)$. Comparing $(i,t)$ with
$(i,t')$ and $(i',t)$ with $(i',t')$ for $t'\in T\setminus\{t\}$, the first coordinate cancels,
so both $\nu$ and $\nu'$ strictly expose $t$ in $T$ for the array $a'$. By Lemma~\ref{lem:combine}
the combination $\nu_0$ strictly exposes $t$ in $T$ and vanishes on $W$ and on $u$, hence on
$W':=W+\R u$. If $T\ne\emptyset$ then $u\notin W$, since $\nu$ vanishes on $W$ but not on $u$;
so $\dim W'=\dim W+1$ and $D\le(r+2)+\dim W=r+1+\dim W'$. The induction hypothesis, applied to
$a'$, $W'$ and $T$, gives $|F_i\cap F_{i'}|\le V(r,m)$.

\smallskip\noindent\emph{Counting.} Lemma~\ref{lem:cs} with $X=[m]^{[r+1]}$, $k=m$ and
$K=V(r,m)$ gives, for $s:=|S|$,
\[
  s^{2}\;\le\;m^{r+1}\bigl(s+m(m-1)V(r,m)\bigr)\;\le\;m^{r+1}s+m^{r+1}m^{2}V(r,m).
\]
A natural number with $s^2\le Us+B$ satisfies $s\le U+\lfloor\sqrt B\rfloor$, so
$s\le V(r+1,m)$.
\end{proof}

The subspace $W$ is the heart of the argument: each step discards one colour, and pays for it by
one more direction which all the relevant functionals annihilate. After $r$ steps a single colour
is left and the admissible functionals form a line.

\begin{corollary}\label{cor:vertexbound}
Let $z\colon[n+1]\times[m]\to\R^{n+1}$ and let $V\subseteq\Grid(n+1,m)$ be convexly independent
for $\Sigma z$. Then $|V|\le V(n,m)$.
\end{corollary}

\begin{proof}
By Lemma~\ref{lem:expose} Theorem~\ref{thm:fibre} applies with $r=n$, $D=n+1$ and $W=0$.
\end{proof}

\subsection{Density of the vertex bound}

\begin{lemma}\label{lem:density}
Let $D\ge1$ and $m:=2^{2^{D}+2}$. For every $r<D$,
\[
  \bigl(V(r,m)+m^{r}\bigr)\cdot2^{2^{D-r}}\;\le\;m^{r+1}.
\]
In particular, for $D=n+1$ and $r=n$, $\ 4\,V(n,m)\le m^{n+1}$.
\end{lemma}

\begin{proof}
For $r=0$ this reads $3\cdot2^{2^D}\le2^{2^D+2}$. For the induction step let $r+1<D$, and put
$p:=m^{r}$, $M:=V(r,m)$ and $T:=2^{2^{D-r-1}}$, so that $T\ge4$ and $T^{2}=2^{2^{D-r}}$. The
induction hypothesis is $(M+p)T^{2}\le pm$, and the claim is
\[
  \bigl(2pm+\lfloor\sqrt{p\,m^{3}M}\rfloor\bigr)\,T\;\le\;pm^{2},
\]
since $V(r+1,m)+m^{r+1}=2pm+\lfloor\sqrt{pm^{3}M}\rfloor$. The hypothesis gives $T^{2}\le m$,
hence $2T\le m$, and $MT^{2}\le p(m-T^{2})$. It therefore suffices that
\[
  p\,m^{3}M\,T^{2}\;\le\;p^{2}m^{3}(m-T^{2})\;\le\;p^{2}m^{2}(m-2T)^{2},
\]
and the second inequality is $m(m-T^2)\le(m-2T)^2$, i.e.\ $mT(T-4)+4T^{2}\ge0$.
\end{proof}

\subsection{Counting small subsets}

\begin{lemma}\label{lem:smallcount}
Let $X$ be a finite set with $M$ elements, let $K\le M$, and let $\mathcal{G}_K$ be the family of
subsets of $X$ with at most $K$ elements. If $4K\le M$, then
\[
  |\mathcal{G}_K|^{4}\cdot2^{3M}\;\le\;3^{4M}.
\]
\end{lemma}

\begin{proof}
Give each $V\subseteq X$ the weight $2^{M-|V|}$. Then $\sum_{V\subseteq X}2^{M-|V|}=3^{M}$, and
every $V\in\mathcal{G}_K$ has weight at least $2^{M-K}$, so $|\mathcal{G}_K|\cdot2^{M-K}\le3^{M}$.
Since $3M\le4(M-K)$, raising to the fourth power gives the claim.
\end{proof}

As in the classical estimate $\sum_{j\le M/4}\binom Mj\le2^{H(1/4)M}$, but with a single weighted
binomial identity and no logarithms.

\section{Assembling the bound}\label{sec:numerics}

Fix $n\ge2$ and set
\[
  D:=n+1,\qquad L:=2^{D}+2,\qquad m:=2^{L},\qquad M:=m^{D}.
\]
By Lemma~\ref{lem:hull}, Corollary~\ref{cor:vertexbound} and Lemma~\ref{lem:density}, every convex
label has a code $(q,V)$ with $|V|\le K:=V(n,m)$ and $4K\le M$. So
$\mathcal{G}=\mathcal{G}_K$ is admissible in Corollary~\ref{cor:product}.

\begin{lemma}[Sign budget]\label{lem:budget}
$20\cdot4L(D+1)^{4}m<M$.
\end{lemma}

\begin{proof}
Dividing by $m$, we must show $80\,L\,(n+2)^{4}<2^{Ln}$. Since $n\ge2$ we have $n+2\le2^{n}$ and
$L=2^{n+1}+2\le2^{n+2}$, so
\[
  80\,L\,(n+2)^{4}\;<\;2^{7}\cdot2^{n+2}\cdot2^{4n}\;=\;2^{5n+9}\;\le\;2^{Ln},
\]
where the last step is $5n+9\le(2^{n+1}+2)n$, true because $2^{n+1}\ge8$.
\end{proof}

\begin{proposition}[Sign factor]\label{prop:signfactor}
$\bigl|\Sigma(D,m)\bigr|^{20}<2^{M}$.
\end{proposition}

\begin{proof}
$L\ge10$, so Corollary~\ref{cor:signpow} applies; conclude by Lemma~\ref{lem:budget}.
\end{proof}

\begin{theorem}\label{thm:labelcount}
For $n\ge2$ and $m=2^{2^{n+1}+2}$, $\ \bigl|\mathcal{L}(n+1,m)\bigr|<2^{m^{n+1}}$.
\end{theorem}

\begin{proof}
Write $P:=|\Sigma(D,m)|$ and $G:=|\mathcal{G}_K|$, so that $|\mathcal{L}|\le PG$ by
Corollary~\ref{cor:product}, $P^{20}<2^{M}$ by Proposition~\ref{prop:signfactor}, and
$G^{4}2^{3M}\le3^{4M}$ by Lemma~\ref{lem:smallcount}. Using $3^{5}\le2^{8}$,
\[
  (PG)^{20}\,2^{15M}\;=\;P^{20}\bigl(G^{4}2^{3M}\bigr)^{5}
  \;<\;2^{M}\,3^{20M}\;\le\;2^{M}\,2^{32M}\;\le\;\bigl(2^{M}\bigr)^{20}2^{15M},
\]
so $PG<2^{M}$.
\end{proof}

The exponent $20$ reflects the balance of the two factors: the small subsets may number up to
$(81/8)^{M/4}\approx2^{0.83M}$, so the sign patterns must stay below $2^{M/6}$ or so.

\begin{proof}[Proof of Theorem~\ref{thm:main} for $n\ge2$]
Theorem~\ref{thm:labelcount} and Corollary~\ref{cor:reduction} give
$\VCn(C)\le m-1=2^{2^{n+1}+2}-1=d(n)$.
\end{proof}

For $n=1$ the sign budget fails ($80\cdot6\cdot3^{4}>2^{6}$), and the next section uses
Theorem~\ref{thm:planar} instead.

\section{The planar bound}\label{sec:planar}

Theorem~\ref{thm:planar} is due to Naszódi and Taschuk \cite[Theorem~3]{naszoditaschuk}, who
prove it for positive homothets of a convex body, i.e.\ of a compact convex set with nonempty
interior; they also construct a convex body in $\R^{3}$ whose translates shatter an infinite
set. The statement for an arbitrary convex $C\subseteq\R^{2}$ follows by a compactness
argument. Suppose that for every $J\subseteq[4]$ there is $y_J$ with $y_J+x_j\in C\iff j\in J$.
The polytope $P:=\conv\{y_J+x_j:j\in J\subseteq[4]\}$ lies in $C$, so it misses the finitely many
points $y_J+x_i$ with $i\notin J$, and some $\varepsilon>0$ is smaller than their distances to
$P$. Then the convex body $P+\varepsilon B$, where $B$ is the closed unit disc, satisfies
$y_J+x_j\in P+\varepsilon B\iff j\in J$, contradicting \cite[Theorem~3]{naszoditaschuk}.

\begin{proof}[Proof of Theorem~\ref{thm:main} for $n=1$]
If $x\colon[1]\times[4]\to\R^{2}$ and $y$ witnessed $\mathrm{VC}_1(C)>3$, the four points
$x(0,j)$ would be shattered by translates of $C$ (they are distinct, since the subsets of $[4]$
are all realised). So $\mathrm{VC}_1(C)\le3\le63=d(1)$.
\end{proof}

\section{The bound \texorpdfstring{$\mathrm{VC}_2\le123$}{VC2 <= 123}}\label{sec:three}

In $\R^{3}$ the same scheme, with sharper bookkeeping, gives Theorem~\ref{thm:three}. Slice the
grid $[m]^{3}$ into the $m$ layers $\{i:i(0)=j\}$. Two layers share at most
\[
  \mu(m):=\Bigl\lfloor\tfrac12\bigl(m+\sqrt{m^{2}+8m^{2}(m-1)}\bigr)\Bigr\rfloor
\]
tails: if they shared three tails, the six points $b_k$, $b_k+u$ ($k<3$) would all be strictly
exposed by functionals annihilating the layer offset $w$, which is impossible because the
functionals annihilating two independent vectors of $\R^{3}$ form a line; a Cauchy--Schwarz count
turns this into the bound $\mu(m)$. A second Cauchy--Schwarz count over the layers shows that a
convexly independent set has at most
\[
  R(m):=\Bigl\lfloor\tfrac12\bigl(m^{2}+\sqrt{m^{4}+4m^{3}(m-1)\mu(m)}\bigr)\Bigr\rfloor
\]
elements. At $m=124$ one has $\mu(124)=2007$ and $R(124)=693785$. With $M=124^{3}=1906624$, the
weighted binomial count with weights $4$ and $7$ bounds the small subsets by
$11^{M}/(4^{R}\,7^{M-R})$, Theorem~\ref{thm:minorsigns} bounds the sign patterns by
$(128e\cdot124^{12}/372)^{1116}<(192\cdot124^{12}+2)^{1118}$, and the inequality
\[
  (192\cdot124^{12}+2)^{1118}\cdot11^{1906624}\;<\;2^{1906624}\cdot4^{693785}\cdot7^{1906624-693785},
\]
between two numbers of about $3.6$ million bits, is checked by exact integer arithmetic. Hence
$|\mathcal{L}(3,124)|<2^{124^{3}}$ and $\mathrm{VC}_2(C)\le123$.

\section{A convex set in \texorpdfstring{$\R^{n+2}$}{R\^{}(n+2)} of infinite
\texorpdfstring{$\VCn$}{VCn} dimension}\label{sec:infinite}

Write points of $\R^{n+2}$ as $(u,h)$ with $u\in\R^{n+1}$ and $h\in\R$, and lift $u$ to the point
$\hat u:=(u,|u|^{2})$ of the paraboloid $h=|u|^{2}$.

\begin{lemma}\label{lem:paraboloid}
Let $T\subseteq\R^{n+1}$ be any set and $q\in\R^{n+1}$. If $\hat q\in\conv\{\hat u:u\in T\}$, then
$q\in T$.
\end{lemma}

\begin{proof}
Consider the affine function $\psi(u,h):=h-2\langle q,u\rangle+|q|^{2}$, which vanishes on the
tangent hyperplane of the paraboloid at $\hat q$. Then $\psi(\hat u)=|u-q|^{2}$, so if
$q\notin T$ then $\psi>0$ on every $\hat u$ with $u\in T$, hence on their convex hull, whereas
$\psi(\hat q)=0$.
\end{proof}

\begin{proof}[Proof of Theorem~\ref{thm:infinite}]
The pairs $(d,S)$ with $d\in\N$ and $S\subseteq[d+1]^{[n]}$ form a countable set; fix an
injection $(d,S)\mapsto\alpha_{d,S}$ of it into $\R$. For such a pair and $i\in[d+1]^{[n]}$ put
$u_{d,S}(i):=(i,\alpha_{d,S})\in\R^{n}\times\R=\R^{n+1}$, and let
\[
  C:=\conv\bigl\{\widehat{u_{d,S}(i)} : d\in\N,\ S\subseteq[d+1]^{[n]},\ i\in S\bigr\}.
\]
Fix $d$. For $k<n$ and $j\in[d+1]$, and for $S\subseteq[d+1]^{[n]}$, put
\[
  x(k,j):=\bigl(j e_k,\,0,\,j^{2}\bigr),
  \qquad
  y(S):=\bigl(0,\,\alpha_{d,S},\,\alpha_{d,S}^{2}\bigr)
\]
in $\R^{n}\times\R\times\R$. Then $y(S)+\sum_{k<n}x\bigl(k,i(k)\bigr)=\widehat{u_{d,S}(i)}$. This
point lies in $C$ if $i\in S$. Conversely, if it lies in $C$, then Lemma~\ref{lem:paraboloid}
gives $u_{d,S}(i)=u_{d',S'}(i')$ for some $d'$, $S'$ and $i'\in S'$; comparing last coordinates
gives $(d',S')=(d,S)$ by injectivity of $\alpha$, and then $i'=i$, so $i\in S$. Thus $x$ and $y$
witness $\VCn(C)>d$.
\end{proof}

For $n=1$ the set $C$ is an unbounded convex set in $\R^{3}$ whose translates have infinite VC
dimension. Unlike the example of \cite{naszoditaschuk}, it is not a convex body, but it realises
the shattering of every finite grid with the same set.

\section{Scope and questions}\label{sec:status}

\subsection{What is not proved}

The case $n=1$ of Question~\ref{q:main}(b) is Theorem~\ref{thm:planar}, due to Naszódi and
Taschuk \cite{naszoditaschuk}. For $n\ge2$, neither (b) nor (c) follows from the results here:
\[
  \mathrm{VC}_2(C)\le2\ \text{ for convex } C\subseteq\R^{3},
  \qquad
  \VCn(C)\le3\ \text{ for convex } C\subseteq\R^{n+1}\ (n\ge2).
\]
The bound $123$ for $n=2$, and the bound $d(n)=2^{2^{n+1}+2}-1$ in general, are far from the
conjectured values. The counting method cannot reach a constant: the grid must have side length
large enough that the number $2^{m^{n+1}}$ of subsets beats the sign patterns, and the side
length $m$ itself is the bound obtained. A proof of (b) or (c) would presumably have to exploit
convexity directly on the shattered configuration, as the proof of Theorem~\ref{thm:planar}
does in the plane.

\subsection{Minkowski sums}

Theorem~\ref{thm:fibre} bounds the number of vertices of a Minkowski sum $P_0+\dots+P_r$ of
point sets $P_k\subseteq\R^{D}$ with $m$ points each by $V(r,m)$, when $D\le r+1$. In this regime
it is a quantitative, elementary relative of the theorem of Sanyal \cite{sanyal} that
$|\vertt(P_1+\cdots+P_r)|<\prod_k|\vertt(P_k)|$ for $r\ge D$, whose proof uses Gale duality and
a topological non-embedding theorem.

\begin{question}
The vertex bound satisfies $V(n,m)=\Theta(m^{n+1-2^{-n}})$ as $m\to\infty$, so convexly
independent sets have density $O(m^{-2^{-n}})$, and this density loss is what forces
$d(n)$ to be doubly exponential. What is the true maximal size of a convexly independent set of
indices of an additive array $[n+1]\times[m]\to\R^{n+1}$? A density of $m^{-c}$ with $c$
independent of $n$ would give a singly exponential bound.
\end{question}

\begin{question}
Can the convexity arguments behind Theorem~\ref{thm:planar} be extended to $\R^{3}$, to prove
$\mathrm{VC}_2(C)\le2$ or at least a small constant bound?
\end{question}

\section{Formalisation}\label{sec:formal}

All the results of this paper are formalised in Lean~4.

\subsection{The development}

The development \cite{repo} is a single Lean~4 project on Lean \lean{v4.33.1} and Mathlib. Its
entry file \lean{VCDimConvexFC.lean} states the finiteness result and
Theorem~\ref{thm:infinite} in terms of an explicit Lean definition of $\VCn$, and
the VC dimension of translates (\lean{HasAddVCDimLE}) is taken from \cite{meanfourier}. The
sources contain no \lean{sorry}, no \lean{admit}, no custom \lean{axiom}, no
\lean{native\_decide} and no \lean{unsafe} declaration; the only large computation is the
certificate of Section~\ref{sec:three}, checked by \lean{decide~+kernel}. The entry file reports
\lean{\#print axioms} for the existence statement, for the bounds $3$, $123$ and $d(n)$, and for
Theorem~\ref{thm:infinite} as
\lean{[propext, Classical.choice, Quot.sound]}.

The results of this paper correspond to the following declarations. Theorem~\ref{thm:planar} is
stated for the VC dimension of the family of translates, in the form \lean{HasAddVCDimLE 3 C},
which implies $\mathrm{VC}_1(C)\le3$.

\begin{center}
\small
\begin{longtable}{@{}p{0.34\linewidth}p{0.60\linewidth}@{}}
\toprule
\textbf{Result here} & \textbf{Lean declaration} \\
\midrule
\endhead
Theorem~\ref{thm:main} & \lean{VCDimConvex.explicit\_bound\_general} \\
Question~\ref{q:main}(a) & \lean{VCDimConvex.fc\_exists\_\allowbreak hasAddVCNDimAtMost\_\allowbreak n\_of\_convex\_rn\_add\_one} \\
Theorem~\ref{thm:planar} & \lean{VCDimConvex.explicit\_bound\_one} \\
Theorem~\ref{thm:three} & \lean{VCDimConvex.explicit\_bound\_123} \\
Theorem~\ref{thm:infinite} & \lean{VCDimConvex.exists\_convex\_\allowbreak forall\_not\_\allowbreak hasAddVCNDimAtMost} \\
Lemma~\ref{lem:paraboloid} & \lean{paraboloidLift\_mem\_convexHull} \\
Proposition~\ref{prop:shatter} & \lean{realizes\_of\_shattering} \\
Corollary~\ref{cor:reduction} & \lean{hasAddVCNDimAtMost\_of\_label\_count\_side} \\
Lemma~\ref{lem:hull} & \lean{exists\_label\_generators} \\
Theorem~\ref{thm:minorhull} & \lean{mem\_indexedHull\_iff\_of\_sameMinorSigns} \\
Corollary~\ref{cor:product} & \lean{card\_convexLabels\_le\_signs\_mul\_generators} \\
Theorem~\ref{thm:minorsigns} (weaker form) & \lean{card\_minorSignPatterns\_le\_direct} \\
Corollary~\ref{cor:signpow} & \lean{card\_minorSignPatterns\_le\_two\_pow\_direct} \\
Lemma~\ref{lem:expose} & \lean{exists\_exposing\_of\_convexIndependentOn} \\
Lemma~\ref{lem:combine} & \lean{exists\_combination\_annihilating} \\
Lemma~\ref{lem:twopoints} & \lean{card\_le\_two\_of\_exposed\_annihilating} \\
Lemma~\ref{lem:cs} & \lean{sum\_card\_sq\_le\_of\_card\_inter\_le} \\
Theorem~\ref{thm:fibre} & \lean{card\_le\_vertexBound\_of\_exposed\_annihilating} \\
Corollary~\ref{cor:vertexbound} & \lean{card\_le\_vertexBound\_of\_convexIndependentOn} \\
Lemma~\ref{lem:density} & \lean{vertexBound\_density}, \lean{four\_mul\_vertexBound\_le} \\
Lemma~\ref{lem:smallcount} & \lean{card\_filter\_card\_le\_mul\_weight\_le},
  \lean{card\_vertexSubsets\_pow\_le} \\
Lemma~\ref{lem:budget} & \lean{sign\_budget} \\
Theorem~\ref{thm:labelcount} & \lean{card\_convexLabels\_lt\_general} \\
Section~\ref{sec:three} & \lean{card\_le\_R\_of\_convexIndependentOn}, \lean{certificate\_124} \\
\bottomrule
\end{longtable}
\end{center}

\subsection{Design choices}

Several arguments were chosen with formalisation in mind. Theorem~\ref{thm:fibre} replaces
arguments based on Gale duality and topological non-embedding theorems by linear algebra and the
Cauchy--Schwarz inequality, and Theorem~\ref{thm:minorhull} avoids any general position
assumption.

Warren's theorem is not in Mathlib, and its usual proofs go through Milnor--Thom bounds or
B\'ezout-type estimates, both heavy to formalise. The formalisation instead proves the weaker
bound $(2kN+1)^p$ on the number of strict sign words of $N$ polynomials of degree $\le k$ in $p$
variables (\lean{card\_strictPatterns\_le\_direct}), with no hypothesis relating $N$ and $p$.
Given a realised sign word, maximise $P^2-\varepsilon\bigl(1+\sum_jx_j^{2kN+2}\bigr)$, where
$P=\prod_if_i$, over a large box intersected with the closed region where the $f_i$ have the
prescribed weak signs. For small $\varepsilon$ the maximiser is an interior point of the region
and of the box, so an unconstrained local maximum, and the maximisers for distinct words are
distinct. The critical equations have the pure-power form $x_j^{2kN+1}=R_j(x)$ with
$\deg R_j\le2kN$, and a linear-algebra argument on the evaluation map shows that such a system
has at most $(2kN+1)^p$ solutions over any field
(\lean{card\_le\_pow\_of\_pure\_power\_relations}). The perturbation of
Corollary~\ref{cor:ternary}, with $\varepsilon$ promoted to a new variable, gives
$(4kN+1)^{p+1}$ ternary sign vectors (\lean{card\_ternaryPatterns\_le\_direct}), and hence the
bound $\bigl|\Sigma(D,m)\bigr|\le\bigl(4D\,\mathcal{N}(D,m)+2\bigr)^{D^2m+2}$ in place of
Theorem~\ref{thm:minorsigns}. This weaker bound still implies Corollary~\ref{cor:signpow}, and it is
the one used in the certificate of Section~\ref{sec:three}.

The formal proof of Theorem~\ref{thm:planar} does not go through \cite{naszoditaschuk}, and
works with $C$ directly. Radon's lemma splits four shattered points into two parts with meeting convex hulls; a part consisting of one
point contradicts convexity at once, and two crossing segments $[x_1,x_3]$, $[x_2,x_4]$ are
excluded by writing $y_{\{1,3\}}+x_2$ or $y_{\{2,4\}}+x_3$ as an explicit convex combination of
the points $y_{\{1,3\}}+x_1$, $y_{\{1,3\}}+x_3$, $y_{\{2,4\}}+x_2$, $y_{\{2,4\}}+x_4$ of $C$ (which
of the two depends on the quadrant of $y_{\{1,3\}}-y_{\{2,4\}}$ in the basis
$x_3-x_1$, $x_4-x_2$).

\subsection{Attribution}

The original formalisation, with the bound $2^{8(n+2)^n}-1$, is due to Kenta Kitamura,
developed with the assistance of large language models, and is released under the Apache~2.0
licence. This paper describes the current state
of the development, in which the fibre induction of Section~\ref{sec:vertices} replaces the
original sparsity argument, and which adds formal proofs of the bounds for $n=1$
(Theorem~\ref{thm:planar}, due to Naszódi and Taschuk \cite{naszoditaschuk}) and $n=2$
(Section~\ref{sec:three}).

\begin{thebibliography}{99}

\bibitem{repo}
K.~Kitamura,
\emph{A Lean proof of a finite additive $\VCn$ bound for convex sets},
\url{https://github.com/KitaKen1/vcdim-convex-finite-bound}.

\bibitem{meanfourier}
Y.~Dillies,
\emph{mean-fourier}, files \lean{Mathlib/Combinatorics/SetFamily/VCDim.lean} and
\lean{Mathlib/Combinatorics/Additive/VCDim.lean},
\url{https://github.com/YaelDillies/mean-fourier}.

\bibitem{naszoditaschuk}
M.~Naszódi and S.~Taschuk,
\emph{On the transversal number and VC-dimension of families of positive homothets of a convex
body},
Discrete Math.\ \textbf{310} (2010), 77--82;
\texttt{arXiv:0907.5223}.

\bibitem{afz}
N.~Alon, J.~Fox and Y.~Zhao,
\emph{Efficient arithmetic regularity and removal lemmas for induced bipartite patterns},
Discrete Anal.\ (2019), Paper No.~3;
\texttt{arXiv:1801.04675}.

\bibitem{chernikovtowsner}
A.~Chernikov and H.~Towsner,
\emph{Hypergraph regularity and higher arity VC-dimension},
\texttt{arXiv:2010.00726}.

\bibitem{cpt}
A.~Chernikov, D.~Palacín and K.~Takeuchi,
\emph{On $n$-dependence},
Notre Dame J.\ Form.\ Log.\ \textbf{60} (2019), 195--214;
\texttt{arXiv:1411.0120}.

\bibitem{gladkova}
V.~Gladkova,
\emph{On a conjecture of Terry and Wolf},
\texttt{arXiv:2411.05612}.

\bibitem{laskowski}
M.~C.~Laskowski,
\emph{Vapnik--Chervonenkis classes of definable sets},
J.\ London Math.\ Soc.\ (2) \textbf{45} (1992), 377--384.

\bibitem{sanyal}
R.~Sanyal,
\emph{Topological obstructions for vertex numbers of Minkowski sums},
J.~Combin.\ Theory Ser.~A \textbf{116} (2009), 168--179;
\texttt{arXiv:math/0702717}.

\bibitem{shelah71}
S.~Shelah,
\emph{Stability, the f.c.p., and superstability; model theoretic properties of formulas in first
order theory},
Ann.\ Math.\ Logic \textbf{3} (1971), 271--362.

\bibitem{shelahstrongly}
S.~Shelah,
\emph{Strongly dependent theories},
Israel J.\ Math.\ \textbf{204} (2014), 1--83;
\texttt{arXiv:math/0504197}.

\bibitem{shelah2dep}
S.~Shelah,
\emph{Definable groups for dependent and $2$-dependent theories},
Sarajevo J.\ Math.\ \textbf{13} (2017), 3--25;
\texttt{arXiv:math/0703045}.

\bibitem{sisask}
O.~Sisask,
\emph{Convolutions of sets with bounded VC-dimension are uniformly continuous},
Discrete Anal.\ (2021), Paper No.~1;
\texttt{arXiv:1802.02836}.

\bibitem{terrywolf}
C.~Terry and J.~Wolf,
\emph{Higher-order generalizations of stability and arithmetic regularity},
\texttt{arXiv:2111.01739}.

\bibitem{terrywolfvc2}
C.~Terry and J.~Wolf,
\emph{The structure of subsets of $\mathbb{F}_p^n$ of bounded $\mathrm{VC}_2$-dimension},
\texttt{arXiv:2510.12867}.

\bibitem{vc}
V.~N.~Vapnik and A.~Ya.~Chervonenkis,
\emph{On the uniform convergence of relative frequencies of events to their probabilities},
Theory Probab.\ Appl.\ \textbf{16} (1971), 264--280.

\bibitem{warren}
H.~E.~Warren,
\emph{Lower bounds for approximation by nonlinear manifolds},
Trans.\ Amer.\ Math.\ Soc.\ \textbf{133} (1968), 167--178.

\end{thebibliography}

\end{document}
